% 本文件由 03_代码/p4_tables.py 自动生成，请勿手工修改。
% 数据源：06_支撑材料/p4_results.json

\begin{table}[htbp]
  \centering
  \caption{4-3 在主计费口径与退款口径下的对照（两口径各自独立标定并回放，模型与执行规则完全相同，只换费用函数）}
  \label{tab:p4-refund}
  \footnotesize
  \begin{tabular}{lrrr}
    \toprule
    指标 & 主口径（不退款、只追加） & 退款口径 & 差异 \\
    \midrule
    合计费用 & 14,196,815.30 & 14,132,052.41 & -64,762.90 \\
    其中 计划费 & 13,266,140.18 & 13,560,753.42 & +294,613.23 \\
    其中 调整费 & 675,425.12 & 352,231.74 & -323,193.39 \\
    其中 紧急费 & 255,249.99 & 219,067.25 & -36,182.74 \\
    上调量 & 564,511.09 & 526,210.68 & -38,300.41 \\
    下调量 & 0.00 & 856,607.01 & +856,607.01 \\
    紧急购电量 & 45,266.69 & 39,357.48 & -5,909.21 \\
    弃电量 & 2,015,140.18 & 1,525,998.13 & -489,142.05 \\
    期末储电量 & 5,912.78 & 5,862.66 & -50.12 \\
    调整提交次数 & 1,002 & 1,002 & --- \\
    \bottomrule
  \end{tabular}
  \par\vspace{2pt}
  \begin{minipage}{.92\textwidth}\footnotesize
    注：退款口径下“少买”不再无条件划算——调减会退回原价，可能诱使计划量虚高、把电量买在未必便宜的时段。因此两种口径的\textbf{标定结果与最终计划都不同}，表中的差异同时包含了参数差异与规则差异，不能读作同一个解的两种计价。这里关心的是\textbf{4-3 相对 4-2 的方向是否依然成立}：只要结论不变，主结论就不依赖这条题面歧义。
  \end{minipage}
\end{table}
